\scanpage{291}

Ceci \deleted{\unsure{posé}} \inserted{\unsure{admis}} on a \inserted{défini} un foncteur
\[
(\pi',\Gamma,I^{\prime !})\longmapsto(\pi,\Gamma,I^{\prime !})
\]
allant de la catégorie des $\Gamma$-groupes ext. à lacets
\inserted{$\pi'$} munis d'une partie $I^{\prime !}$ de $I'=I(\pi')$ stable
\inserted{\deleted{sur lesquels $\Gamma$ opère} telle que $\unsure{s'}\in I'\Rightarrow\Gamma_{\unsure{s'}}\ne(1)$,}
par $\Gamma$, dans la catégorie des $\Gamma$-groupes extérieurs à lacets, munis d'une partie $I^{\prime !}$ de
$\Phi(\pi,\Gamma)$. La conjecture est que le foncteur
(lui-$\overline{\mathrm m}$ un peu conjectural, sauf si on se réduit aux situations de départ réalisables, d'où situations d'arrivée également réalisables) est une équivalence de catégories, i.e. $1^\circ)$ pl. fidèle et $2^\circ)$ \unsure{ess.} surjectif. \deleted{Le point}

Le point $2^\circ)$ est clair, quand on se borne de part et d'autre à des situations réalisables. Mais même dans ce cas, le point $1^\circ)$ n'est pas prouvé.
\deleted{\unsure{Il se peut que}} \inserted{Je pense} que
\deleted{\unsure{les propriétés} \illegible{2} \unsure{sur le} \illegible{2} \unsure{fidélité}}
\inserted{\uline{si} $I'$ \uline{est réduit à} \uline{un élément}, alors les}
\deleted{\[
\operatorname{Autext}(\pi',I^{\prime !})\longrightarrow
\operatorname{Autext}(\unsure{\pi)\Gamma|\pi,\Gamma},\unsure{I^{\prime !}|I''}
\]}
propriétés du foncteur ``forage d'un trou'' permettront de l'établir aisément. On voit aussi que pour l'établir, on est

\scanpage{292}
\mspage{147}

ramené à établir la bijectivité des $(*)$, et le \unsure{th.} d'équivalence dans le cas où $\Gamma$ est transitif sur $I^{\prime !}$.

Pour résumer, le Programme d'attaque du Pb de réalisabilité d'un $\Gamma$-groupe extérieur à lacets serait le suivant

I) Etablir la bijectivité des $(*)$ dans le cas général

(II) \unsure{Établir} la pleine fidélité du foncteur précédent (en se \unsure{ramenant} au besoin au cas où $I^{\prime !}$ est \uline{une} orbite (singulière) de $\Gamma$ dans $I'=I(\pi')$.

(III) Dans le cas où $\Phi(\pi,\Gamma)=\varnothing$ i.e. l'extension $\unsure{\mathcal E}$ n'a pas d'élément d'ordre fini $\ne1$, prouver que $\unsure{\mathcal E}$, muni \inserted{de l'ens.} des $\unsure{\mathcal E}$-classes de conjugaison des centralisateurs des $L_i$ $(i\in I(\pi))$ dans $\unsure{\mathcal E}$, est un groupe à lacets.

D'ailleurs, pour disposer des propriétés préliminaires indispensables des ens. $\Phi(\pi,\Gamma)$, il faudrait commencer par réaliser ce programme pour les

\scanpage{293}

groupes cycliques (\unsure{opérant} de façon directe), et procéder dans ce cas par dévissage.

Si le th. de classification complet (comme équivalence des 2-catégories \inserted{--des opérations} topologiques et des opérations \unsure{isotopiques|homologiques} --) était vrai, les pts I, II, III de ce programme \unsure{devraient} l'être aussi, et ce serait donc un bon programme d'approche. Les pts I et II devraient être valables, dès que le 2-foncteur entre 2-catégories serait 2-fidèle (pas néc. \deleted{pas} 3-fidèle), \deleted{et la condition néc.} \inserted{\unsure{du moins pour des}} situations réalisables. On dirait donc qu'une action \inserted{\unsure{ext.}} de $\Gamma$ sur un $\pi'$, muni d'un $I^{\prime !}$, est \uline{admissible}, si \deleted{toute}
\[
\operatorname{Aut\,ext}(\pi',\Gamma,I^{\prime !})\xrightarrow{\sim}\operatorname{Aut}(\pi,\Gamma,I^{\prime !})
\]
et \unsure{si} de plus tte \deleted{\unsure{bonne}} autre action de $\Gamma$ sur $\pi'$, \deleted{induisant la $\overline{\mathrm m}$ action sur $I^{\prime !}$ et} la $\overline{\mathrm m}$ action sur $\pi$ et la $\overline{\mathrm m}$ application $I^{\prime !}\unsure{\longrightarrow|{-}}\Phi(\pi,\Gamma)$, est conjuguée de l'action originelle

\scanpage{294}
\mspage{148}

Ceci posé, la cond. \unsure{néc.} et suff. de \unsure{réalisabilité} d'une action ext. \inserted{fidèle} de $\Gamma$ sur un $\pi$ serait alors que cette action se \unsure{remonte} \inserted{(par ``\unsure{forage} de trous'' \deleted{\unsure{pour $I^{\prime !}$}} \inserted{pour $I^{\prime !}=\Phi(\pi,\Gamma)$})} en une action \inserted{\uline{admissible},} de $\Gamma$ sur un $(\pi',I^{\prime !})$, \deleted{\illegible{1}} (par définition $\overline{\mathrm m}$, ce \unsure{remontement} serait \deleted{que l'ex. corr. de $\Gamma$ \unsure{par} \unsure{des} $\pi'$,} alors \unsure{ess.} unique) et que de plus l'ext. correspondante $\unsure{\mathcal E}'$ de $\Gamma$ par $\pi'$ soit un groupe à lacets.

25. Groupes de Teichmüller \uline{spéciaux}.

Revenons aux notations du n° 19, On va définir un sous-groupe $SA_{g,\nu}^{\unsure{\deleted{!}|!}}$ de $A_{g,\nu}^{!}$
\[
SA_{g,\nu}^{\unsure{\deleted{!}|!}}\unsure{\{} =
\{u\in A_g\mid u\text{ induise l'identité sur un voisinage de }S_{g,\nu}\}
\]
\inserted{(ens. des \unsure{autom.} de $U_{g,\nu}$ qui)}
i.e \deleted{$u$} \unsure{induisent|induit} l'identité \unsure{ds|dans} le complémentaire d'un compact.

Contrairement aux autres ss-groupes de $A_g$ considérés jusqu'à présent, celui-ci n'est pas un ss-groupe fermé. \uline{NB} Dans le cas où on travaille sur des surfaces compactes

\scanpage{295}

à bord au lieu de surfaces compactes (sans bord) ``\unsure{trouées|troués}'', il y aurait lieu de prendre le groupe des autom. qui \unsure{induisent} l'identité sur le bord.

$SA_{g,\nu}$ est un ss-gr \unsure{invariant} de $A_{g,\nu}$, le \unsure{quot} $A_{g,\nu}/SA_{g,\nu}$ étant \unsure{isom.} au groupe des \uline{germes} d'\unsure{autom.} de $X_g$ au voisinage de $\unsure{S_{g,\nu}|SA_{g,\nu}}$, ou encore le groupe des germes \inserted{à l'infini} d'\unsure{autom.} de $\unsure{U_{\nu,g}|U_{g,\nu}}$ (dans les complémentaires de compacts\ldots)

\deleted{Considérons} \inserted{On aura} évidemment, puisque
\[
SA_{g,\nu}\subset A_{g,\nu}
\]
\[
(SA_{g,\nu})^{0}\subset \unsure{SA_{g,\nu}^{0}|A_{g,\nu}^{0}},
\]
\unsure{posons}
\[
S\Gamma_{g,\nu}^{!}=SA_{g,\nu}/SA_{g,\nu}^{0},
\]
on aura un homomorphisme canonique
\[
S\Gamma_{g,\nu}^{!}\longrightarrow\Gamma_{g,\nu}^{!+}.
\]
qu'on va interpréter de façon algébrique, en termes de l'interprétation de $\Gamma_{g,\nu}^{!+}$ \inserted{$=\Gamma^{!}$} comme

\scanpage{296}
\mspage{149}

le groupe des autom. ext. du groupe à lacets $\pi_{g,\nu}$, \unsure{induisant} l'identité sur $S_{g,\nu}=I(\pi_{g,\nu})$.

\deleted{Pour tt \illegible{1}} \deleted{Soit} Pour $i\in I=S_{g,\nu}$, choisissons un \deleted{$\unsure{\ell}$} $L_i$ de classe $i$ dans $\pi_{g,\nu}$ --- ce qui revient à se donner un ``pt'' \unsure{du|de} $\unsure{B}_{D_i^{*}}$ \inserted{($\simeq$ un rev. universel $\widetilde D_i^{*}$)} \deleted{\unsure{sur}} et un isom \unsure{entre} son image \unsure{dans} $\unsure{B}_{U=U_{g,\nu}}$ avec l'\deleted{\unsure{élément}} ``point'' $s=s_{g,\nu}$ de référence, qui servait pour définir $\pi_{g,\nu}$ comme $\operatorname{Aut}_{\unsure{B}_U}(s_{\deleted{\unsure{g,\nu}}})\simeq\pi_1(\unsure{B}_U,s)$.

Ceci dit, si $u$ est un autom. de $\pi=\pi_{g,\nu}$, le fait qu'il respecte (\unsure{strict.}, \inserted{en induisant l'identité sur $S_{g,\nu}$)} la str. à lacets, s'exprime par l'existence d'une famille d'élts $g_i\in\unsure{\pi}$, \unsure{tels} que
\[
u(\ell)=\operatorname{int}g_i^{-1}\bigl(\deleted{\unsure{\chi(u)}}\ell^{\deleted{\unsure{\chi(u)}}\inserted{\alpha}}\bigr)
\qquad(i\in I,\unsure{\forall}\ell\in L_i)
\]
\margin{où $\alpha=\chi(u)$}
--- \unsure{lesquels|les quels} $g_i$ sont déterminés par $u$ \deleted{\illegible{1}} modulo multiplication à droite par un $\lambda_i\in L_i$. Si on a de $\overline{\mathrm m}$ $v,(h_i)$, alors \deleted{sur $L_i$} Pour $\ell\in L_i$ \unsure{on a}
\[
\begin{aligned}
(vu)(\ell)
&=v\bigl(\underbrace{u(\ell)}_{\inserted{\operatorname{int}(g_i^{-1})\ell^{\alpha}}}\bigr)
=v\bigl(\operatorname{int}(g_i^{-1})\bigr)v(\ell)^{\alpha}\\
&=\operatorname{int}\bigl(v(g_i^{-1})\bigr)\operatorname{int}\bigl(\unsure{h_i^{-1}}\bigr)\cdot(\ell^{\alpha\beta})\\
&=\operatorname{int}\bigl(v(g_i^{-1})\cdot h_i^{-1}\bigr)\cdot\ell^{\alpha\beta}
=\operatorname{int}\bigl((h_i\cdot v(g_i))^{-1}\bigr)\ell^{\alpha\beta}
\end{aligned}
\]
\margin{NB $\alpha=\chi(u)$\quad $\beta=\chi(v)$}
donc $vu$ est compatible avec le système des \deleted{$\unsure{u|v}(g_i)$} $h_i v(g_i)$ \unsure{Posons}

\scanpage{297}

\[
(v,(h_i))\cdot(u,(g_i))=(vu,h_i v(g_i)\deleted{\unsure{h_i}}),
\]
\margin{$S$}

On trouve alors une structure de groupe sur l'ens \inserted{$S\unsure{\mathcal E}^{!}$} des $(u,(g_i))$, sous-groupe du produit semi-direct de $\unsure{\mathcal E}^{!}=\mathrm{Autext}_{\mathrm{lac}}(\pi,I)$ \inserted{id. sur $I$} par $\pi^I$, sur lequel $\unsure{\mathcal E}^{!}$ opère de façon évidente. L'hom naturel
\[
S\unsure{\mathcal E}^{!}\longrightarrow\unsure{\mathcal E}^{!}
\]
est surjectif, et son noyau est can. isom à $\prod_{i\in I} L_i\simeq \unsure{T|\mathbb T}^I$, où $\unsure{T|\mathbb T}$ est le module des orientations.
\deleted{Je voudrais définir un \unsure{hom} canonique
\[
S\Gamma_{g,\nu}^{!}\longrightarrow S^{\prime}\Gamma^{!}
\]
compatible avec les \unsure{hom} des $\Gamma^{!}$, ce qui définit}
\inserted{d'où} une structure d'extension \deleted{(\unsure{centrale})}
\margin{où $\unsure{\mathcal E}$ opère sur $\unsure{T|\mathbb T}^I$ via son action sur $\unsure{T|\mathbb T}$}
\[
1\longrightarrow \unsure{T|\mathbb T}^I\longrightarrow S\unsure{\mathcal E}^{!}\longrightarrow\unsure{\mathcal E}^{!}-1
\]
\deleted{
\[
1\unsure{\longrightarrow}\unsure{T|\mathbb T}_g^{S_{g,\nu}}\unsure{\longrightarrow}S\Gamma_{g,\nu}^{!}\unsure{\longrightarrow}\Gamma_{g,\nu}^{!}-1
\]
\unsure{sur} $S\Gamma_{g,\nu}^{!}$ (où $\unsure{T|\mathbb T}_g\,(\unsure{\simeq}\unsure{\mathbb Z})$ est le module d'orientations). Il faut donc définir

Notons que le noyau de $S\Gamma_{g,\nu}^{!}\longrightarrow\Gamma_{g,\nu}^{!}$ s'identifie à
$SA_{g,\nu}^{\unsure{\deleted{!}|!}}\cap A_{g,\nu}^{0}/(SA_{g,\nu}^{\unsure{\deleted{!}|!}})^{0}$, et
}

\scanpage{298}
\mspage{150}

\deleted{
\[
SA_{g,\nu}\longrightarrow S^{\prime}\Gamma^{!}
\]
en utilisant le choix, pour $\unsure{\forall|\mathrm{tt}}\ i\inserted{\in I=S_{g,\nu}}$ d'un rev. universel $\widetilde D_i^{*}$ d'un voisinage $D_i^{*}$ de $s_i$ (\unsure{qu'on} prendra toujours : un disque épointé) et d'un \unsure{iso} entre le rev. universel $\unsure{\widetilde U(i)|U(i)}$ de $U=U_{g,\nu}$ qu'il définit, et le rev. \unsure{univ.} de référence $\widetilde U$ (pointé au dessus de $s=s_{g,\nu}$). Cet \unsure{autom}
}

Je voudrais \unsure{interpréter} cette extension comme l'image inverse d'une extension \inserted{can} de $\Gamma^{!}$ par $\unsure{T|\mathbb T}^I$ (canonique en tous cas, une fois choisis les $\widetilde D_i^{*}$). Pour ceci, notons que si $\unsure{\ddot u|\dot u}$ est une \unsure{autoéquivalence|autoéquiv.} \deleted{\unsure{dans}} de la situation $\unsure{B}_{D^{*}}\longrightarrow\unsure{B}_U$ \inserted{induisant l'identité sur $I$} on peut considérer les $\dot u(\widetilde D_i^{*})$, et on peut considérer les $\widetilde{\dot u}$ au dessus de $\dot u$, savoir les syst. $(\dot u,\gamma_i)$, où pour \unsure{tt|tout} $i\in I$, $\gamma_i:\widetilde D_i^{*}\longrightarrow\dot u(\widetilde D_i^{*})$ [\unsure{déterminé} \unsure{mod|par} \unsure{élt|élémt} de $\unsure{T|\mathbb T}$]. \deleted{\unsure{L'ens des}} Le groupe des classes d'isomorphie d'\unsure{automs|autom.} de $(\unsure{B}_{D^{*}}\unsure{\longrightarrow}\unsure{B}_U,(\widetilde D_i^{*}))$, ou si on préfère, des classes d'isom. d'\unsure{automs|autom.} de
\[
(\unsure{B}_I\longrightarrow\unsure{B}_{D^{*}}\unsure{\longrightarrow}\unsure{B}_U)
\]

\scanpage{299}

induisant l'identité sur $I$, est donc une extension de $\Gamma^{!}$ par $\unsure{T|\mathbb T}^I$, soit $S\Gamma^{!}$. On va définir un hom. de suites exactes
\diagram{\begin{tikzcd}[ampersand replacement=\&,column sep=2em,row sep=2.2em]
1 \arrow[r,no head,to path={(\tikztostart.east) -- ([xshift=-2pt]\tikztotarget.west) -- ++(0,-0.07) \tikztonodes}]
\& T^I \arrow[r,no head] \arrow[d,equal]
\& S\unsure{\mathcal E}^{!}
 \arrow[r,no head,to path={(\tikztostart.east) -- ([xshift=-2pt]\tikztotarget.west) -- ++(0,-0.07) \tikztonodes}]
 \arrow[d,no head]
\& \deleted{S}\unsure{\mathcal E}^{!}
 \arrow[r,no head,to path={(\tikztostart.east) -- ([xshift=-2pt]\tikztotarget.west) -- ++(0,-0.10) \tikztonodes}]
 \arrow[d]
\& 1
\\
\deleted{\unsure{\illegible{1}\,\Gamma}}
 \arrow[r,no head,to path={(\tikztostart.east) -- ([xshift=-2pt]\tikztotarget.west) -- ++(-0.03,-0.06) \tikztonodes}]
\& T^I
 \arrow[r,no head,to path={(\tikztostart.east) -- ([xshift=-2pt]\tikztotarget.west) -- ++(-0.03,-0.06) \tikztonodes}]
\& S\Gamma^{!} \arrow[r,no head]
\& \Gamma^{!} \arrow[r,no head]
\& 1
\end{tikzcd}}
qui prouve que $S\unsure{\mathcal E|\varepsilon}^{!}$ est bien l'image inverse de \deleted{\illegible{1}} l'extension \inserted{$S\Gamma^{!}$} de $\Gamma^{!}$ par $\unsure{T|\mathbb T}^I$.

Pour cela, on note que $\deleted{S}\unsure{\mathcal E|\varepsilon}^{!}$ est le groupe des classes d'isom. \unsure{d'autoéqui.|de autoéqui.} de
\diagram{\begin{tikzcd}[ampersand replacement=\&]
\unsure{B}_{D^*} \arrow[d] \& {} \\
\unsure{B}_U \& \text{pt.}s \arrow[l]
\end{tikzcd}}
et $S\unsure{\mathcal E|\varepsilon}^{!}$ celui des classes d'isom des \deleted{\illegible{1}} autoéqui. de
\diagram{\begin{tikzcd}[ampersand replacement=\&,column sep=2.4em,row sep=1.7em]
B_{D^*} \arrow[dd]
\& {} \& B_I \arrow[ll] \arrow[dd]
\\
{} \& {} \& {}
\\
B_U
\& {} \arrow[ul,bend right=20,shorten <=5pt,shorten >=5pt,"{\begin{gathered}\text{isom}\\[-2pt]\text{de}\\[-2pt]\text{comm.}\end{gathered}}"'{xshift=1.3em,yshift=-1.8em}]
\& \text{pt.}\,s \arrow[ll]
\end{tikzcd}}
qui s'envoie de façon naturelle dans celui des classes d'iso. des autoéqui. des \unsure{2-diagrammes}
\diagram{\begin{tikzcd}[ampersand replacement=\&]
\unsure{B}_I \arrow[r] \& \unsure{B}_{D^*} \arrow[r] \& \unsure{B}_U
\end{tikzcd}}.

\scanpage{300}
\mspage{151}

En fait, dans tout ceci il n'y avait aucunement lieu de se borner aux \deleted{\unsure{extérieurs}} \inserted{\unsure{autom.}} de $\pi$ (ext. ou totaux) induisant l'identité sur $I$. On trouve de toutes façons des extensions $S\unsure{\mathcal E|\varepsilon}$ de $\unsure{\mathcal E|\varepsilon}$ par $\unsure{T|\mathbb T}^I$, $S\Gamma$ de $\Gamma$ par $\unsure{T|\mathbb T}^I$, et un hom d'extensions
\diagram{\begin{tikzcd}[ampersand replacement=\&]
{} \& {} \& 1 \arrow[d] \& 1 \arrow[d] \& {} \\
{} \& {} \& \pi \arrow[r,"\unsure{\sim}"] \arrow[d] \& \pi \& {} \\
1 \arrow[r] \& \unsure{T|\mathbb T}^I \arrow[r] \arrow[d,equal] \& S\unsure{\mathcal E|\varepsilon} \arrow[r] \arrow[d] \& \unsure{\mathcal E|\varepsilon} \arrow[r] \arrow[d] \& 1 \\
1 \arrow[r] \& \unsure{T|\mathbb T}^I \arrow[r] \& S\Gamma \arrow[r] \arrow[d] \& \Gamma \arrow[r,no head] \arrow[d] \& 1 \\
{} \& {} \& 1 \& 1 \& {}
\end{tikzcd}}
de façon que l'on pt considérer $S\unsure{\mathcal E|\varepsilon}$ comme ext. de $\Gamma$ par $\unsure{T|\mathbb T}^I\times\pi$
\[
1\longrightarrow \unsure{T|\mathbb T}^I\times\pi\longrightarrow S\unsure{\mathcal E|\varepsilon}\longrightarrow\Gamma\longrightarrow1
\]
dont \inserted{\unsure{l'ext}} $\unsure{\mathcal E|\varepsilon}$ resp. $S\Gamma$ se déduit en divisant par $\unsure{T|\mathbb T}^{\unsure{I|S}}$ resp. par $\pi$.

\unsure{\deleted{Revenant}|Revenant} maintenant au cas-type \deleted{$U_{g,\nu}$}
$\unsure{U}_{g,\nu}=X_{g,\nu}\setminus S_{g,\nu}$, on \unsure{prendra|prend} \unsure{les|des} notations
\margin{\deleted{\unsure{$S\mathcal E_{g,\nu}$}}}
\[
\unsure{\mathcal E|\varepsilon}_{g,\nu}
=\mathrm{Aut}\deleted{\unsure{\mathrm{ext}}}_{\mathrm{lac}}
(\unsure{\pi|\Pi}_{g,\nu})
\]
($\stackrel{\mathrm{can}}{\simeq}\Gamma_{g,\nu+1}$ si $(g,\nu)\ne(0,0),(0,1)$ i.e. \deleted{\unsure{$\pi_{g,\nu}\ne0$}}) \illegible{1}

et
\[
S\Gamma_{g,\nu}=S\mathrm{Autext}_{\mathrm{lac}}(\unsure{\pi|\Pi}_{g,\nu})
\]
(extension de $\Gamma_{g,\nu}$ par $\unsure{T|\mathbb T}^I$)

et $S\unsure{\mathcal E|\varepsilon}_{g,\nu}$
($\simeq S\Gamma'_{g,\nu+1}/\unsure{T|\mathbb T|\Gamma}_{\unsure{s_\nu|g,\nu}}$
\inserted{\deleted{\illegible{1}} \unsure{où} $S\Gamma'_{g,\nu+1}$ désigne le}
sous-groupe de $S\Gamma'_{g,\nu+1}$ des \unsure{élé-ts} qui fixent \unsure{le} dernier élément $s_\nu$ \ldots)
